% SPDX-License-Identifier: MIT
\chapter{Voraussetzungen, Build und Installation}\label{chap:install}
\ChapterSummary{Sie richten eine reproduzierbare C++20-Buildumgebung ein,
starten die Demonstration, führen Tests aus und binden das installierte
Toolkit in ein eigenes CMake-Projekt ein. Unterschiede zwischen
Ein-Konfigurations- und Visual-Studio-Builds werden ausdrücklich erklärt.}

\section{Voraussetzungen und Quellen}
Sie benötigen einen C++20-fähigen Compiler, CMake und ein
Buildsystem für Ihre Plattform. Das Haupt-Toolkit ist bewusst ohne
obligatorische SDL-Abhängigkeit nutzbar. Unter Linux und macOS ist
ein interaktives ANSI/VT-Terminal für die sichtbare Terminal-Demo
sinnvoll; unter Windows wird eine geeignete Konsole bzw. ein
Terminal mit VT-Unterstützung benötigt.

Der kanonische Quellstand liegt unter
\url{https://github.com/Robin-Goerlach/SASD-UI-Toolkit}.
Klonen Sie das Repository oder arbeiten Sie in Ihrem vorhandenen
Arbeitsverzeichnis. Bauen Sie \emph{außerhalb} des Quellbaums in
einem eigenen Build-Verzeichnis. Dadurch bleiben Quellstand, Cache,
Tests und Varianten getrennt.

\section{Erster Build unter Linux oder macOS}
\begin{lstlisting}[language=bash,caption={Konfigurieren, bauen und testen}]
cmake -S . -B build \
  -DCMAKE_BUILD_TYPE=Debug \
  -DSASD_UI_BUILD_TESTS=ON \
  -DSASD_UI_BUILD_EXAMPLES=ON
cmake --build build --parallel
ctest --test-dir build --output-on-failure
\end{lstlisting}
\term{Konfigurieren} prüft Compiler, Features und Abhängigkeiten.
\term{Bauen} erstellt Bibliotheken, Beispiele und Tests.
\term{Testen} führt die registrierten CTest-Fälle aus.
Schlägt schon die Konfiguration fehl, hilft ein Blick in die erste
aussagekräftige CMake-Fehlermeldung; das spätere Weiterbauen mit
einem inkonsistenten Cache ist meist keine gute Idee.

\section{Visual Studio 2022 unter Windows}
\begin{lstlisting}[language={},caption={Multi-Config-Build mit MSVC}]
cmake -S . -B build-msvc -G "Visual Studio 17 2022" -A x64 ^
  -DSASD_UI_BUILD_TESTS=ON ^
  -DSASD_UI_BUILD_EXAMPLES=ON
cmake --build build-msvc --config Debug --parallel
ctest --test-dir build-msvc -C Debug --output-on-failure
\end{lstlisting}
Die Fortsetzungszeichen \texttt{\char94} sind für \emph{cmd.exe}
gedacht; unter PowerShell verwenden Sie entweder eine einzige Zeile
oder deren Backtick-Zeilenfortsetzung. Bei Visual-Studio-Generatoren
muss die Konfiguration \filepath{Debug} oder \filepath{Release}
sowohl beim Build als auch bei CTest angegeben werden.

\section{Die Terminal-Demo starten}
Unter einem üblichhen Linux-Ein-Konfigurations-Build:
\begin{lstlisting}[language=bash]
./build/examples/sasd_ui_terminal_demo
\end{lstlisting}
Unter Windows ist der Pfad beispielsweise
\filepath{build-msvc/examples/Debug/sasd_ui_terminal_demo.exe}.
Eine wirkliche Terminal-Sitzung ist erforderlich; die Anwendung
lässt sich nicht sinnvoll als nichtinteraktiver Batch-Prozess testen.

Erkunden Sie die Oberfläche: Mit \textbf{Tab} und
\textbf{Shift+Tab} wechseln Sie den Fokus. \textbf{Enter}
oder \textbf{Leertaste} aktiviert geeignete Buttons.
\textbf{F1} zeigt die Demo-Hilfe, \textbf{F10} öffnet das Menü.
Die Demo unterstützt unter anderem Texteingabe, Checkbox,
Radiobuttons und eine ComboBox. Beenden ist über den Exit-Button,
Escape oder den vorgesehenen Demo-Shortcut möglich.

\Important{Terminalgroesse, Mausprotokoll und
Steuersequenzen hängen vom Hostterminal ab. Die \term{TerminalSession}
verwaltet die Zustandsänderungen per RAII und stellt sie bei normalem
Verlassen und beim Abwickeln von C++-Ausnahmen wieder her.
Eine harte Prozessbeendigung kann diese Garantie nicht bieten.}

\section{Release-Build und Installation}
\begin{lstlisting}[language=bash,caption={Installation in ein gewähltes Präfix}]
cmake -S . -B build-release \
  -DCMAKE_BUILD_TYPE=Release \
  -DSASD_UI_BUILD_TESTS=OFF \
  -DSASD_UI_BUILD_EXAMPLES=OFF
cmake --build build-release --parallel
cmake --install build-release --prefix /opt/sasd-ui
\end{lstlisting}
Verwenden Sie unter Windows bei Bedarf zusätzlich
\filepath{--config Release}. Ein eigenständiges Installationspräfix
ist für Tests und mehrere Versionen meist sicherer als eine
ungefragte systemweite Installation.

\section{Das Toolkit als CMake-Paket verwenden}
\begin{lstlisting}[language={},caption={Auszug aus CMakeLists.txt}]
cmake_minimum_required(VERSION 3.20)
project(my_ui_app LANGUAGES CXX)
set(CMAKE_CXX_STANDARD 20)
set(CMAKE_CXX_STANDARD_REQUIRED ON)

find_package(SASDUIToolkit 0.2 CONFIG REQUIRED)
add_executable(my_ui_app main.cpp)
target_link_libraries(my_ui_app PRIVATE SASD::UI::Terminal)
\end{lstlisting}
Bei einem abweichenden Installationsort geben Sie ihn über
\filepath{CMAKE_PREFIX_PATH} an, zum Beispiel:
\begin{lstlisting}[language=bash]
cmake -S . -B build -DCMAKE_PREFIX_PATH=/opt/sasd-ui
\end{lstlisting}
Die installierten Targets haben unterschiedliche Aufgaben:
\begin{itemize}
\item \texttt{SASD::UI}: backendneutraler Core.
\item \texttt{SASD::UI::Terminal}: Terminal-Backend.
\item \texttt{SASD::UI::Rendered}: SDL-unabhängige Rendered-Schicht.
\end{itemize}
Verlinken Sie nur die Schicht, die Ihre Anwendung tatsächlich benötigt.
Die Versionsnummer der Paketanforderung ist an den genutzten
Entwicklungs- bzw. Release-Stand anzupassen: Das veröffentlichte
v0.1.0-Tag bildet den älteren Terminal-Baseline-Stand ab.

\section{Optional: SDL3 bewusst einschalten}
Für die experimentelle Desktop-Demo existieren die Optionen
\filepath{SASD_UI_BUILD_SDL3_ADAPTER=ON} und
\filepath{SASD_UI_FETCH_SDL3=ON}. Hinzu kommt bei den
entsprechenden Tests ein Fontpfad über
\filepath{SASD_UI_SDL3_TEST_FONT}.
Das normale Toolkit braucht weder SDL3 noch SDL\_ttf.
Der Adapter wird gegenwärtig \emph{nicht} wie Core, Terminal und
Rendered als stabiler externer CMake-Installationsvertrag angeboten.
Weitere Details finden Sie in \filepath{README.md} und
\filepath{docs/en/M3_RENDERED_DESKTOP_SMOKE_TEST.md}.

\section{Den erfolgreichen Build überprüfen}
Bevor Sie ein eigenes Projekt aufsetzen, prüfen Sie vier voneinander
unabhängige Punkte: CMake hat die gewünschten Targets erzeugt;
der Compiler und Linker haben keine Fehler gemeldet; CTest
hat die vorgesehenen Testfälle wirklich ausgeführt; und die
Terminal-Demo lässt sich interaktiv bedienen. Ein erfolgreiches
Kompilieren allein bestätigt noch nicht den tatsächlichen
Eingabe- und Darstellungspfad.

\begin{lstlisting}[language=bash,caption={Build- und Testumfang nachvollziehen}]
cmake --build build --target help
ctest --test-dir build -N
ctest --test-dir build --output-on-failure
\end{lstlisting}
Die erste Zeile listet die vom gewählten Generator angebotenen
Targets; die zweite zeigt die von CTest registrierten Testfälle,
ohne sie auszuführen. Fehlen erwartete Tests, prüfen Sie
\filepath{SASD_UI_BUILD_TESTS} und den CMake-Cache. Für
Fehleranalysen sollten Sie die erste fehlgeschlagene
Testausgabe mitsamt Compiler- und Betriebssystemversion
festhalten, statt ungezielt Build-Dateien zu verändern.

\Important{Ein von CTest bestätigter Headless-Test beweist nicht
automatisch, dass Ihr konkretes Terminal alle benötigten
Tasten, Mausprotokolle oder Unicode-Breitenkonventionen
unterstützt. Testen Sie die Referenz-Demo zusätzlich in der
später vorgesehenen Zielumgebung.}

\PracticalTip{Bei Problemen zuerst die Demo und die Tests aus dem
gleichen Checkout bauen. Wenn diese funktionieren, liegt ein Fehler
bei Ihrem Konsumentenprojekt häufig an Link-Targets,
Paketpfad oder einer gemischten Debug/Release-Konfiguration.}
